% Draft text for Section 4.2 (Datasets): query strings for the MS-only comparison against the
% repositioning oracle (reproduction of TeraTools Figure 5, with Orbit modes added).
% Provenance: /vast/blangme2/terams-fig (bench/make_patterns.sh, bench/sample_patterns.py,
% patterns/*.fa and *.fa.tsv); index/pattern split from Nate Brown's
% /vast/blangme2/nbrown/ms_lcp/data/<collection>/pattern/<collection>.5%.fa.
% Collection sizes (N, n, r, M, m) are from TeraTools Table S8.

\paragraph{Held-out query strings.}
To compare the MS-only algorithms with each other and with the repositioning oracle, we used the three collections and the index/query split from the TeraTools matching statistics benchmark~\cite{teratools}.
Each collection consists of the forward and reverse-complement strands of its genomes; 5\% of these sequences were withheld as queries and the remaining 95\% were indexed.
Both strands of a withheld sequence are withheld together, so no query occurs in the index as its own reverse complement.
% verified for mtb152 (16 sequences = both strands of 8 genomes); for chr19.1000 (100 = 2 x 50) and human100 inferred from the split, verify.
The queries are therefore complete genomic sequences from individuals related to, but absent from, the index, so the matches are long and repositioning is rare (0.01\% to 0.07\% of query characters in TeraTools' measurements).
The withheld sequences contain only A, C, G and T.
\begin{itemize}
  \item \textbf{mtb152}: 152 \textit{Mycobacterium tuberculosis} genomes. We used all 16 withheld sequences, both strands of 8 complete genomes (70.6\,Mbp; 4.39 to 4.43\,Mbp each).
  \item \textbf{chr19.1000}: 1{,}000 chromosome~19 haplotypes from 1000 Genomes Phase~3, generated with \texttt{bcftools consensus}. Of the 100 withheld sequences (55.0\,Gbp), we used 19 whole sequences (1.03\,Gbp; 48.8 to 62.8\,Mbp each).
  \item \textbf{human100}: 100 human haplotype assemblies, including those of HPRC release~1. Of the 3{,}864 withheld sequences (63.8\,Gbp), we used 97 whole sequences (1.09\,Gbp; 1.4\,kbp to 179\,Mbp each).
\end{itemize}
For chr19.1000 and human100, we chose whole withheld sequences in random order (fixed seed 20261006) until their total length reached 1\,Gbp, and kept them in their original order.
Sequences were not split or trimmed.
The oracle records the outcome of every repositioning step (the target BWT interval and the longest common extension) for a given index and query set, and replays it at query time, so it does no candidate scan and no LCE computation.
Its time and operation count are a lower bound for any method that shares the same LF-step machinery.

% Bib entry to add to refs.bib (verify fields before use; Table S8/S9 and Figure 5 are in this paper):
% @article{teratools,
%   title={{RLBWT}-Based Index Construction in Compressed Space for Terabase-Scale Pangenome Analysis},
%   author={Brown, Nathaniel K. and others},
%   year={2026}}
